\documentclass{amsart}
% For dealing with references we use the comment environment
\usepackage{verbatim}
\newenvironment{reference}{\comment}{\endcomment}
%\newenvironment{reference}{}{}
\newenvironment{slogan}{\comment}{\endcomment}
\newenvironment{history}{\comment}{\endcomment}

% For commutative diagrams we use Xy-pic
\usepackage[all]{xy}

% We use 2cell for 2-commutative diagrams.
\xyoption{2cell}
\UseAllTwocells

% We use multicol for the list of chapters between chapters
\usepackage{multicol}

% This is generally recommended for better output
\usepackage{lmodern}
\usepackage[T1]{fontenc}

% For cross-file-references

% Package for hypertext links:
\usepackage{hyperref}

% For any local file, say "hello.tex" you want to link to please

% Theorem environments.
%
\theoremstyle{plain}
\newtheorem{theorem}[subsection]{Theorem}
\newtheorem{proposition}[subsection]{Proposition}
\newtheorem{lemma}[subsection]{Lemma}

\theoremstyle{definition}
\newtheorem{definition}[subsection]{Definition}
\newtheorem{example}[subsection]{Example}
\newtheorem{exercise}[subsection]{Exercise}
\newtheorem{situation}[subsection]{Situation}

\theoremstyle{remark}
\newtheorem{remark}[subsection]{Remark}
\newtheorem{remarks}[subsection]{Remarks}

\numberwithin{equation}{subsection}

% Macros
%
\def\lim{\mathop{\mathrm{lim}}\nolimits}
\def\colim{\mathop{\mathrm{colim}}\nolimits}
\def\Spec{\mathop{\mathrm{Spec}}}
\def\Hom{\mathop{\mathrm{Hom}}\nolimits}
\def\Ext{\mathop{\mathrm{Ext}}\nolimits}
\def\SheafHom{\mathop{\mathcal{H}\!\mathit{om}}\nolimits}
\def\SheafExt{\mathop{\mathcal{E}\!\mathit{xt}}\nolimits}
\def\Sch{\mathit{Sch}}
\def\Mor{\mathop{\mathrm{Mor}}\nolimits}
\def\Ob{\mathop{\mathrm{Ob}}\nolimits}
\def\Sh{\mathop{\mathit{Sh}}\nolimits}
\def\NL{\mathop{N\!L}\nolimits}
\def\CH{\mathop{\mathrm{CH}}\nolimits}
\def\proetale{{pro\text{-}\acute{e}tale}}
\def\etale{{\acute{e}tale}}
\def\QCoh{\mathit{QCoh}}
\def\Ker{\mathop{\mathrm{Ker}}}
\def\Im{\mathop{\mathrm{Im}}}
\def\Coker{\mathop{\mathrm{Coker}}}
\def\Coim{\mathop{\mathrm{Coim}}}

% Boxtimes
%
\DeclareMathSymbol{\boxtimes}{\mathbin}{AMSa}{"02}

%
% Macros for moduli stacks/spaces
%
\def\QCohstack{\mathcal{QC}\!\mathit{oh}}
\def\Cohstack{\mathcal{C}\!\mathit{oh}}
\def\Spacesstack{\mathcal{S}\!\mathit{paces}}
\def\Quotfunctor{\mathrm{Quot}}
\def\Hilbfunctor{\mathrm{Hilb}}
\def\Curvesstack{\mathcal{C}\!\mathit{urves}}
\def\Polarizedstack{\mathcal{P}\!\mathit{olarized}}
\def\Complexesstack{\mathcal{C}\!\mathit{omplexes}}
% \Pic is the operator that assigns to X its picard group, usage \Pic(X)
% \Picardstack_{X/B} denotes the Picard stack of X over B
% \Picardfunctor_{X/B} denotes the Picard functor of X over B
\def\Pic{\mathop{\mathrm{Pic}}\nolimits}
\def\Picardstack{\mathcal{P}\!\mathit{ic}}
\def\Picardfunctor{\mathrm{Pic}}
\def\Deformationcategory{\mathcal{D}\!\mathit{ef}}

\setlength{\emergencystretch}{3em}
\begin{document}
\title[Stacks Project, Tag 0GVK]{430 --- Bass freeness theorem for countably generated infinite-rank projectives}
\maketitle
\noindent Copyright (C) 2005--2025 Johan de Jong. Stacks Project authors. Source: \href{https://stacks.math.columbia.edu/tag/0GVK}{Tag 0GVK}.

\noindent Editorial difficulty note (not author text). Difficulty H3: The proof constructs an increasing sequence of finite stably free direct summands capturing all generators. It then extracts an infinite free summand by successive good-element choices and applies an Eilenberg swindle to absorb the complementary countably generated projective module..

\noindent Editorial provenance note (not author text). History: excerpt from revision \texttt{a0444\allowbreak 6e57e\allowbreak c1fbc\allowbreak 252a8\allowbreak 71afc\allowbreak ec775\allowbreak 2fb28\allowbreak 07b14}, more-algebra.tex, lines 39152--39209. Reference presentation was adapted; original-author context and core support blocks were retained or added. The mathematical wording is unchanged. Licensed under GNU FDL 1.2 or later; no invariant sections or cover texts. See ../licenses/COPYING. Context blocks reproduce author definitions and supporting results; they are not additional counted problems.

\begin{theorem}
\label{theorem-countable-free}
\begin{reference}
Commutative case of \href{https://stacks.math.columbia.edu/bibliography}{Bibliography: Bass (Theorem 3.1)}
\end{reference}
Let $R$ be a ring with Jacobson radical $J$ such that $R/J$ is Noetherian.
Let $P$ be a countably generated projective $R$-module such that
$P_\mathfrak m$ has infinite rank for all maximal ideals
$\mathfrak m$ of $R$. Then $P$ is free.
\end{theorem}

\begin{proof}
We first prove that $P$ is a countable direct sum of finite stably
free modules. Let $x_1, x_2, \ldots$ be a countable set of generators for $P$.
We inductively construct finite stably free direct summands
$F_1, F_2, \ldots$ of $P$ such that for all $n$ we have
that $F_1 \oplus \ldots \oplus F_n$ is a direct summand of $P$
which contains $x_1, \ldots, x_n$. Namely, given $F_1, \ldots, F_n$
with the desired properties, write
$$
P = F_1 \oplus \ldots \oplus F_n \oplus P'
$$
and let $s \in P'$ be the image of $x_{n + 1}$.
By Lemma \href{https://stacks.math.columbia.edu/tag/0GVJ}{lemma element in free summand (Tag 0GVJ)}
we can find a finite stably free direct summand $F_{n + 1} \subset P'$
containing $s$. Then $P = \bigoplus_{i = 1}^{\infty} F_i$.

\medskip\noindent
Assume that $P$ is an infinite direct sum
$P = \bigoplus_{i = 1}^{\infty} F_i$ of nonzero
finite stably free modules. The stable freeness
of the modules $F_i$ will be used in the following
manner: the rank of each $F_i$ is constant (and positive). Hence
we see that $P_\mathfrak m$ is free
of countably infinite rank for each maximal ideal
$\mathfrak m$ of $R$. By Lemma \href{https://stacks.math.columbia.edu/tag/0GVI}{lemma trick to find good element (Tag 0GVI)}
applied with $s = 0$ and $M = P$, we can find
a $t_1 \in P$ such that $Rt_1$ is a free direct summand of $P$.
Then $t_1$ is contained in $F_1 \oplus \ldots \oplus F_{n_1}$
for some $n_1 > n_0 = 0$.
The same reasoning applied to $\bigoplus_{n > n_1} F_n$
produces an $n_1 < n_2$ and $t_2 \in F_{n_1 + 1} \oplus \ldots \oplus F_{n_2}$
which generates a free direct summand.
Continuing in this fashion we obtain a free direct summand
$$
\bigoplus\nolimits_{i \geq 1} t_i :
\bigoplus\nolimits_{i \geq 1} R
\longrightarrow
\bigoplus\nolimits_{i \geq 1}
\bigoplus\nolimits_{n_i \geq n > n_{i - 1}} F_n = P
$$
of infinite rank. Thus we see that $P \cong Q \oplus F$ for some free
$R$-module $F$ of countable rank. Since $Q$ is countably
generated it follows that $Q \oplus Q' \cong F$ for
some module $Q'$. Then the Eilenberg swindle
(Lemma \href{https://stacks.math.columbia.edu/tag/0GVF}{lemma eilenberg swindle (Tag 0GVF)}) implies that
$Q \oplus F \cong F$ and $P$ is free.
\end{proof}
\end{document}
